\section{Height separation and a weak minimizing enclosure}
\label{sec:height}

Throughout this section the data satisfy Definition~\ref{def:prepared}.  In particular, $K$ has compact
support, the dominant energy inequality is strict, and every component of
$S$ has strictly negative future expansion.  Write $a$ for the full-cut
infimum of these prepared data.  Fix $0<\epsilon<1$.  We use the smooth
solution on the filled manifold $\mathcal M_N$ supplied by
Theorem~\ref{thm:elliptic:filled}, with the notation of
Proposition~\ref{prop:scalar:identity}.  Thus
\begin{equation}\label{eq:height:parameters}
 \ell=e^{-N\epsilon},\qquad \eta=\ell^{3/2},\qquad h=\eta f,
 \qquad \widehat g=e^{2t}(g+l^2df\otimes df).
\end{equation}
On the filling, $g$ denotes its smooth extension.  All comparisons with
the original cut infimum will take place in the original exterior.

We first state precisely what is needed from the analytic construction.
After increasing a polynomial $P_N=P(N)$, its estimates give
\begin{equation}\label{eq:height:quantitative-input}
 -\epsilon<t\le Z\le P_N,\qquad
 \ell\le l\le e,\qquad
 |f|+|df|_g\le e^{P_N},\qquad |h|\le C_0.
\end{equation}
Moreover, $|h|\le Cr^{-b_1}$ on the original end, with $C$ independent
of $N$.  Local geometry of the background filling is uniformly bounded;
its added product lengths and volume are $O(1+N)$.  In particular it has
uniform coordinate balls of some radius $r_0>0$, with uniform local
isoperimetric constants.  For each fixed $N$, the end estimates are
\begin{equation}\label{eq:height:end-input}
 t=O_j(r^{2-n}),\qquad \Delta_g t=O(r^{-n-\beta}),
 \qquad f\text{ and its derivatives decay exponentially},
\end{equation}
where $\beta>0$.  Constants in Equation~\eqref{eq:height:end-input}
may depend arbitrarily on fixed $N$.  Only
Equation~\eqref{eq:height:quantitative-input} and the uniform height
barrier will be used in estimates required to be polynomial in $N$.

The height separates the filling from infinity: it remains uniformly
negative on the filling and tends uniformly to zero on the original
end. We use this separation to place a large constant conformal
factor around the filling. The cutoff vanishes on a distant tail,
so it leaves the ADM energy unchanged. Its transition derivatives
are small by weighted coercivity. We then minimize perimeter subject
to a distance obstacle inside the plateau. Contact with that obstacle
would give positive volume densities on both sides of the frontier;
the plateau would turn the resulting perimeter into an area larger
than a distant competitor. Detachment places the exterior and a
neighborhood of its frontier in the region of controlled scalar
curvature. The construction and its estimates are at fixed prepared
data and fixed $\epsilon$.

\subsection{A negative height on the whole filling}

\begin{lemma}[The cap and collar barrier]\label{lem:height:ramp}
The product lengths in the filling can be chosen $O(1+N)$ so that, for
some $b_2>0$ independent of $N$, every sufficiently large $N$ satisfies
\begin{equation}\label{eq:height:negative-core}
 h<-b_2
 \quad\text{on the whole added filling and a neighborhood of }S.
\end{equation}
The conclusion concerns the final solution with the original prepared
$K$; no analogous assertion is required along the existence homotopy.
\end{lemma}

\begin{proof}
Use the signed transverse coordinate $s$ increasing from each product
neck toward $S$ and the original end.  The fixed transition from product
geometry to an original collar has been chosen so that its leaves have
\begin{equation}\label{eq:height:leaf-expansion}
 H_s+\tr_{TS_s}K\le-4\delta
\end{equation}
for some fixed $\delta>0$.  Shrink $\delta$ when necessary.  On the cap
and the product part $K=-L_0g$, where $L_0$ is a fixed large constant.

It is useful to write the trace equation directly in terms of $h$.  At a
point where a test function has the same gradient as $h$, put
$a_0=\eta/l$ and let $\chi$ have its value for the solution there.  The
operator applied to the test function $v_*$ is
\begin{equation}\label{eq:height:test-operator}
 \mathcal Q(v_*)=
 \tr_{A_\chi}K+
 \frac{\tr_{A_\chi}\Hess_g v_*}
 {\sqrt{a_0^2+|dv_*|_g^2}}.
\end{equation}
The equation is $\mathcal Q(h)=h$.  The coefficients in
Equation~\eqref{eq:height:test-operator} are evaluated at the solution;
we do not replace its value of $t$ by one computed from the barrier.
Nevertheless,
\begin{equation}\label{eq:height:axial-smallness}
 0<\chi\le d=\frac{a_0^2}{a_0^2+|dh|_g^2},
 \qquad \frac\eta e\le a_0\le\frac\eta\ell=\ell^{1/2}.
\end{equation}
These inequalities suffice for a comparison at a positive maximum of
$h-v_*$.

Set $v_*=-\delta$ on the cap.  On a product neck let
$q_N=v_*'\ge0$ start from zero, increase approximately exponentially,
and end at a small fixed number $q_*>0$.  One explicit construction is
to take an exponential $q_*e^{s-L_N}$ where its size lies between a
constant multiple of $\eta$ and $q_*$, smooth its end to the constant
$q_*$ on an interval of fixed length, and multiply its initial part by
a fixed smooth cutoff on an interval where it is $O(\eta)$.  Choose
\[
 L_N=\log(q_*/\eta)+O(1)=O(1+N).
\]
The cutoff can be flat at its initial endpoint.  This construction gives
\begin{equation}\label{eq:height:ramp-estimates}
 |q_N'|\le C(q_N+\eta),\qquad
 \int q_N\,ds\le Cq_*+C\eta,
\end{equation}
with $C$ independent of $N$.  Thus $v_*$ glues smoothly to the constant
on the cap.  The same construction is made separately on every neck.
The cap value is common to all components.

The product Hessian has only an axial component.  Equations
\eqref{eq:height:axial-smallness}--\eqref{eq:height:ramp-estimates}
give, even where the slope vanishes,
\[
 \frac{\chi|v_*''|}{\sqrt{a_0^2+|v_*'|^2}}
 \le C\frac{q_N+\eta}{\sqrt{(\eta/e)^2+q_N^2}}\le C'.
\]
Consequently $\mathcal Q(v_*)\le-kL_0+C'$ on this part.  Choose
$L_0$ large enough that this is less than $-2\delta$.  Since
$v_*\ge-\delta$, this is a strict supersolution inequality
$\mathcal Q(v_*)<v_*$.  It also holds on the constant cap, where the
trace of $K$ is $-nL_0$.

Continue with the fixed slope $q_*$ across the fixed transition and up
to $S$.  Choose $q_*$ so small that all the increase up to $S$, including
Equation~\eqref{eq:height:ramp-estimates}, is less than $\delta/4$.
Beyond $S$, still inside the fixed strictly trapped collar, increase the
slope smoothly until $v_*>C_0+1$ at the outer edge of that collar.  This
requires only a fixed, possibly large, slope and fixed derivatives;
they do not depend on $N$.  Keep $v_*'\ge q_*$ on the entire transition
and collar.  For such a function of signed distance,
\[
 \mathcal Q(v_*)=
 \tr_{TS_s}K+\chi K(\partial_s,\partial_s)
 +\frac{v_*'}{\sqrt{a_0^2+(v_*')^2}}H_s
 +\frac{\chi v_*''}{\sqrt{a_0^2+(v_*')^2}}.
\]
Equation~\eqref{eq:height:axial-smallness} shows that this converges
uniformly to the left side of Equation~\eqref{eq:height:leaf-expansion}.
For example, the absolute error is bounded by
$C\ell/q_*^2+C\ell\sup|v_*''|/q_*^3$, with fixed $C$.
For large $N$ it is therefore less than $\delta$.  As
$v_*\ge-\delta$, again $\mathcal Q(v_*)<v_*$.

Apply the maximum test on the compact region consisting of the cap,
necks, transitions, and these outer collars.  On its outer boundary
$h-v_*<0$.  At a positive interior maximum, the gradients agree and
$\Hess h\le\Hess v_*$.  Positive definiteness of $A_\chi$ would give
\[
 h=\mathcal Q(h)\le\mathcal Q(v_*)<v_*<h,
\]
a contradiction.  Thus $h\le v_*$.  Up to $S$ the barrier is less
than $-3\delta/4$.  Continuity of the fixed collar barrier gives the
same strict negative bound on a small neighborhood of $S$.  Taking,
for example, $b_2=\delta/2$ proves the assertion.
\end{proof}

\subsection{Increasing area without losing exterior scalar curvature}

Choose $0<b_3<b_2/4$ so small that
\begin{equation}\label{eq:height:strict-margin}
 \mu-|J|_g-\rho-b_3|\tau|\ge\tfrac12\rho>0
 \quad\text{on the original exterior}.
\end{equation}
This follows from the strict prepared inequality and the choice of
$\rho$ in the scalar construction: outside the compact support of
$\tau$ no smallness condition on $b_3$ is needed, and on that compact
set choose $b_3|\tau|\le\rho/2$.  Let $\zeta\in C^\infty(\R)$ be
nonnegative, equal to zero for $|z|\le b_3/4$, equal to one for
$|z|\ge b_3/2$, and with values in $[0,1]$.  Define
\begin{equation}\label{eq:height:plateau-metric}
 A_N=(\ell/\eta)^{1/2}=\ell^{-1/4},\qquad
 \psi_N= A_N\zeta(h),\qquad
 \widetilde g_N=e^{2\psi_N}\widehat g.
\end{equation}
The height barrier on the end puts $\{|h|\ge b_3/4\}$ in a fixed
original coordinate radius, together with the added filling.  Hence
$\psi_N$ vanishes identically near infinity.  It changes neither the
ADM energy nor the asymptotic estimates of $\widehat g$.

\begin{lemma}[Curvature on the small-height exterior]
\label{lem:height:curvature}
For all sufficiently large $N$,
\[
 R_{\widetilde g_N}\ge0
 \quad\text{on }\Omega\cap\{|h|\le b_3\}.
\]
On the original exterior the metric satisfies
$\widetilde g_N\ge e^{-2\epsilon}g$.  Its ADM energy obeys
\begin{equation}\label{eq:height:energy-upper}
 E(\widetilde g_N)\le E+P_N(\ell+\ell^{1/2}).
\end{equation}
\end{lemma}

\begin{proof}
The solution supplied by Theorem~\ref{thm:elliptic:filled} has
$t>-\epsilon$ and the parameters of the scalar construction. Thus
Equation~\eqref{eq:scalar:solution-curvature} applies without a change
of data or normalization. On $|h|\le b_3$, its constraint and trace
terms have lower bound $\rho/2$ by
Equation~\eqref{eq:height:strict-margin}; the remaining terms are
nonnegative and retain the positive contribution
$(1-\delta_0)\mathcal T$.

For the new conformal factor we use
Equation~\eqref{eq:scalar:height-estimates}, with
$\eta=\ell^{3/2}$. Its constants are polynomial in $N$ for these
fixed prepared data. The factor $\sqrt d$ in the height Laplacian
cancels the axial loss in weighted coercivity, as shown immediately
after that equation. No arbitrary fixed-$N$ estimate for unweighted
second derivatives enters this bound.

The conformal scalar-curvature law in dimension $n=k+1$ reads
\begin{align}\label{eq:height:conformal-curvature}
 \tfrac12e^{2(t+\psi_N)}R_{\widetilde g_N}
 =\tfrac12e^{2t}R_{\widehat g}
 -k e^{2t}\Delta_{\widehat g}\psi_N
 -\tfrac{k(k-1)}2e^{2t}|d\psi_N|_{\widehat g}^2.
\end{align}
The chain rule, Equation~\eqref{eq:scalar:height-estimates}, and the fixed derivative bounds of
$\zeta$ imply
\begin{align*}
 |e^{2t}\Delta_{\widehat g}\psi_N|
 &\le P_N\ell^{1/4}(1+\sqrt{\mathcal T})+C\ell^{3/4},\\
 e^{2t}|d\psi_N|_{\widehat g}^2&\le C\ell^{1/2}.
\end{align*}
On the support of these errors $b_3/4\le|h|\le b_3/2$.  This region
is disjoint from the filling by Lemma~\ref{lem:height:ramp}, and is
contained in a fixed compact portion of the original exterior.  There
$\rho\ge c_*>0$.  Young's inequality bounds the error containing
$\sqrt{\mathcal T}$ by
$(1-\delta_0)\mathcal T/2+P_N\ell^{1/2}$.  All remaining errors
tend to zero.  Equation~\eqref{eq:scalar:solution-curvature} therefore proves
nonnegative scalar curvature there for large $N$.  Where $\zeta$ is
constant, Equation~\eqref{eq:height:conformal-curvature} has no error
terms and the same conclusion follows directly.

The floor $t>-\epsilon$, the nonnegative graph term, and
$\psi_N\ge0$ give the metric comparison.  Finally $\psi_N=0$ near
infinity, so Equation~\eqref{eq:height:energy-upper} is exactly
Proposition~\ref{prop:scalar:flux}.  In particular its sign is the
upper energy estimate required for the Riemannian inequality.
\end{proof}

\subsection{Perimeter minimization and detachment from a distance obstacle}

We use full perimeter in the filled manifold.  For a measurable set
$F$ and an open set $U$, write
$P_\gamma(F;U)=|D\mathbf1_F|_\gamma(U)$ for perimeter in a smooth
metric $\gamma$, and $P_\gamma(F)=P_\gamma(F;\mathcal M_N)$.
The frontier is the support of this perimeter measure.  For a set of
finite perimeter its integration boundary is the reduced boundary
$\partial^*F$, and
$P_\gamma(F;U)=\mathcal H^k_\gamma(\partial^*F\cap U)$.
These conventions distinguish full perimeter from a relative frontier
inside $\Omega$.

The final exterior and an ambient neighborhood of its frontier must
lie in the original small-height region where
Lemma~\ref{lem:height:curvature} applies. We therefore enclose the
whole set where $|h|\ge b_3$, which contains the filling, and use a
thin distance neighborhood of that set as the obstacle. Its distance
collar will supply the filled-side density if a minimizing frontier
touches it. Put
\begin{equation}\label{eq:height:obstacle}
 C_N^h=\{|h|\ge b_3\},\qquad
 O_N=\{x:\dist_g(x,C_N^h)\le s_N\},
 \qquad s_N=e^{-Q_N},
\end{equation}
where $Q_N$ is a sufficiently large polynomial.  It can be chosen so
that $s_N<r_0/10$ and
\begin{equation}\label{eq:height:plateau-neighborhood}
 3s_N\sup|dh|_g<b_3/2.
\end{equation}
Indeed $|dh|\le\eta e^{P_N}$ by
Equation~\eqref{eq:height:quantitative-input}.  The entire
$2s_N$-neighborhood of $O_N$ then lies in $\{|h|>b_3/2\}$, where
$\widetilde g_N=e^{2A_N}\widehat g$.  The compact obstacle $O_N$
contains the whole added filling and a neighborhood of $S$.
No smoothness of $\partial O_N$ is asserted or needed. Outside $O_N$
we are in $\operatorname{int}\Omega\cap\{|h|<b_3\}$; detachment from
$O_N$ will provide the required ambient neighborhood as well.

\begin{lemma}[Existence with a compact support bound]
\label{lem:height:BV-existence}
There is a bounded finite-perimeter set $F_N\supset O_N$, modulo null
sets, minimizing $P_{\widetilde g_N}$ among all bounded sets containing
$O_N$.  One may choose it to contain every other minimizer modulo null
sets.  Its perimeter satisfies
\begin{equation}\label{eq:height:competitor-bound}
 P_{\widetilde g_N}(F_N)\le e^{P_N}.
\end{equation}
\end{lemma}

\begin{proof}
At fixed $N$ the end is smooth and asymptotically Euclidean with the
differentiated falloff in Equation~\eqref{eq:height:end-input}.
Consequently sufficiently large coordinate spheres have positive
mean curvature for $\widetilde g_N$ toward infinity.
The required radius may depend arbitrarily on $N$.

Here is the clipping argument that gives the compact support needed
for the direct method.  All objects in this paragraph use
$\widetilde g_N$.  On the large end let
$Y=\nabla r/|\nabla r|$, so $|Y|=1$ and
$\divg Y=H_r>0$.  For a bounded finite-perimeter competitor $F$ and
almost every large $R$, Gauss--Green on $F\cap\{r>R\}$ gives
\[
 \int_{F\cap\{r>R\}}H_r\,dV
 =\int_{\partial^*F\cap\{r>R\}}\langle Y,\nu_F\rangle\,dA
 -\mathcal H^k(F^{(1)}\cap S_R).
\]
The spherical normal here is $-Y$.  Boundedness of $F$ eliminates an
outer flux; alternatively first cut off $Y$ beyond its support.  The
first boundary integral is at most $P(F;\{r>R\})$.  The BV
truncation formula therefore yields the stronger clipping inequality
\begin{equation}\label{eq:height:clipping}
 P_{\widetilde g_N}(F\cap\{r<R\})
 \le P_{\widetilde g_N}(F)
 -\int_{F\cap\{r>R\}}H_r\,dV_{\widetilde g_N}.
\end{equation}
The Gauss--Green and truncation formulas hold at almost every radius
by the trace and slicing theorems for BV functions \cite{AFP}.
The notation $\{r<R\}$ here includes the entire compact
core.  In particular it contains $O_N$ once $R$ is sufficiently large.

Choose a fixed interval of such radii beyond the obstacle.  Applying
Equation~\eqref{eq:height:clipping} to each term of a minimizing
sequence, with an almost-everywhere admissible radius in that interval,
confines the sequence to one compact set.  BV compactness gives
$L^1$ convergence along a subsequence, and lower semicontinuity gives
a minimizer.  Containment of $O_N$ is preserved almost everywhere.
The minimum is the unrestricted bounded-competitor minimum, so local
variations are not subject to an artificial outer constraint.

Choose a fixed original coordinate sphere $S_{R_*}$ so far out that
$|h|<b_3/4$ on and beyond it for every $N$.  It encloses $O_N$ when
$N$ is large.  On it $\psi_N=0$, and
Equation~\eqref{eq:height:quantitative-input} implies
$\widehat g\le e^{P_N}g$.  Its filled side is a competitor; increasing
$P_N$ gives Equation~\eqref{eq:height:competitor-bound}.  The fixed
sphere need not lie in the region used for clipping: the latter was
needed only for compactness at fixed $N$.

All minimizers have a common compact support bound: if a minimizer had
positive volume beyond an almost-everywhere admissible clipping radius,
the strictly positive last integral in
Equation~\eqref{eq:height:clipping} would contradict minimality.
Among minimizers, maximize the enclosed volume.  Compactness,
lower semicontinuity, and convergence of volume show that this maximum
is attained.  Perimeter submodularity implies
\[
 P(F\cup G)+P(F\cap G)\le P(F)+P(G).
\]
For two minimizers both sets on the left contain $O_N$, so equality
holds and their union is a minimizer.  Maximal volume therefore implies
$G\subset F_N$ modulo null sets for every minimizer $G$.
\end{proof}

\begin{lemma}[Two densities and detachment]\label{lem:height:detachment}
For all sufficiently large $N$,
\begin{equation}\label{eq:height:detached}
 \supp|D\mathbf1_{F_N}|\cap O_N=\varnothing.
\end{equation}
\end{lemma}

\begin{proof}
The two different density arguments are illustrated in
Figure~\ref{fig:height:two-densities}.  Suppose $p$ belongs to this
intersection.  It cannot lie in the
interior of $O_N$, because $F_N$ contains $O_N$ almost everywhere.
Thus $\dist_g(p,C_N^h)=s_N$.  The complete background metric has a
minimizing geodesic from some $z\in C_N^h$ to $p$.  In particular
$B^g_{s_N}(z)\subset O_N$.  The midpoint of this geodesic is the
center of a ball of radius $s_N/4$ contained in both $O_N$ and
$B^g_{s_N}(p)$.  Uniform background ball-volume bounds give
\begin{equation}\label{eq:height:filled-density}
 \Vol_g(F_N\cap B^g_{s_N}(p))\ge c s_N^n.
\end{equation}
This is the filled-side density, obtained from the distance obstacle.

We derive the other density from minimization.  On $B^g_{s_N}(p)$
the high conformal multiplier is constant.  Remove it and write
$\gamma=\widehat g$.  Equation~\eqref{eq:height:quantitative-input}
gives, on this ball and every $k$-plane,
\begin{equation}\label{eq:height:pointwise-comparison}
 e^{-P_N}g\le\gamma\le e^{P_N}g,
 \qquad e^{-P_N}dA_g\le dA_\gamma\le e^{P_N}dA_g,
\end{equation}
after changing the polynomial.  No derivative of $\gamma$ appears
in this comparison.

Put $U=\mathcal M_N\setminus F_N$ and
$V(r)=\Vol_g(U\cap B^g_r(p))$.  Since $p$ is in the support of
perimeter, $V(r)>0$ for every $r>0$: otherwise $\mathbf1_{F_N}$
would be constant almost everywhere on some neighborhood of $p$.
For almost every $0<r<s_N$, filling the ball is an admissible
competitor.  Cancel the unchanged exterior perimeter and the common
constant factor $e^{kA_N}$.  The perimeter truncation formula gives
\[
 P_\gamma(U;B^g_r(p))
 \le\mathcal H^k_\gamma(U^{(1)}\cap\partial B^g_r(p)).
\]
Here $U^{(1)}$ is the measure-theoretic interior; the displayed trace
agrees with it for almost every radius.  The coarea formula in the
smooth background metric and Equation~\eqref{eq:height:pointwise-comparison}
then yield
\begin{equation}\label{eq:height:one-sided-perimeter}
 P_g(U;B^g_r(p))\le e^{P_N}V'(r)
 \quad\text{for almost every }r\in(0,s_N).
\end{equation}

The absolute isoperimetric inequality on a uniformly controlled
background ball, applied to the zero extension of
$\mathbf1_{U\cap B^g_r(p)}$, includes its spherical trace:
\[
 V(r)^{(n-1)/n}
 \le C\bigl(P_g(U;B^g_r(p))+V'(r)\bigr)
 \le e^{P_N}V'(r).
\]
Local Euclidean isoperimetry in a uniform coordinate ball proves this
form; the smooth background metric changes its constant by a fixed
factor.  Since $V$ is absolutely continuous and positive for $r>0$,
its $n$th root satisfies
\[
 (V^{1/n})'(r)\ge n^{-1}e^{-P_N}
 \quad\text{almost everywhere on }(0,s_N).
\]
Integrate from $a>0$ to $r$ and let $a\downarrow0$.  The result is
\begin{equation}\label{eq:height:unfilled-density}
 V(r)\ge e^{-P_N}r^n\qquad(0<r\le s_N),
\end{equation}
with another polynomial.  This proves the unfilled-side density from
a one-sided comparison, without assuming that the obstacle contact is
a regular point or that the deformed metric has controlled derivatives.

Relative isoperimetry in $B^g_{s_N}(p)$, together with Equations
\eqref{eq:height:filled-density} and \eqref{eq:height:unfilled-density},
now gives
\[
 P_g(F_N;B^g_{s_N}(p))\ge e^{-P_N}s_N^{n-1}.
\]
Absorb $s_N^{n-1}=e^{-(n-1)Q_N}$ into the polynomial exponential,
use Equation~\eqref{eq:height:pointwise-comparison}, and restore the
constant multiplier.  We obtain
\begin{equation}\label{eq:height:contact-cost}
 P_{\widetilde g_N}(F_N)
 \ge\exp(kA_N-P_N)
 =\exp\bigl(k e^{N\epsilon/4}-P_N\bigr).
\end{equation}
For fixed $\epsilon>0$, $e^{N\epsilon/4}$ dominates every polynomial
in $N$.  Equation~\eqref{eq:height:contact-cost} contradicts
Equation~\eqref{eq:height:competitor-bound} for large $N$.
\end{proof}

\begin{figure}[t]
\centering
\begin{tikzpicture}[scale=0.92,>=Stealth,
  every node/.style={font=\small}]
  \fill[blue!9] (0,0) circle[radius=1.65];
  \draw[blue!55!black,thick] (0,0) circle[radius=1.65];
  \draw[orange!65!black,thick,dashed] (1.65,0) circle[radius=1.65];
  \fill[blue!40] (0.825,0) circle[radius=0.4125];
  \draw[blue!60!black] (0.825,0) circle[radius=0.4125];
  \draw[black!65] (0,0)--(1.65,0);
  \node[font=\scriptsize,fill=blue!9,inner sep=1pt] at (0.50,0.78) {$d_g(z,p)=s$};
  \fill (0,0) circle[radius=0.035] node[below left] {$z$};
  \fill (0.825,0) circle[radius=0.035];
  \fill (1.65,0) circle[radius=0.045] node[below right] {$p$};
  \node[blue!55!black,align=center] at (-0.55,1.90)
    {$B_s^g(z)\subset O_N$};
  \node[orange!65!black] at (2.27,1.90) {$B_s^g(p)$};
  \draw[->,blue!60!black] (0.78,-1.78)--(0.825,-0.46);
  \node[align=center,anchor=north] at (0.78,-1.90)
    {a ball of radius $s/4$\\inside both balls};
  \node[align=left,anchor=west] at (4.0,0.4)
    {Obstacle inclusion gives\\
     $\Vol_g(F_N\cap B_s^g(p))\ge c s^n$.\\[7pt]
     Filling-ball comparison gives\\
     $\Vol_g(F_N^c\cap B_s^g(p))\ge e^{-P_N}s^n$.};
\end{tikzpicture}
\caption{The two densities at a hypothetical contact point
$p\in\supp|D\mathbf1_{F_N}|\cap\partial O_N$, with $z\in C_N^h$ chosen so that $d_g(z,p)=s=s_N$.
The circles depict background metric balls schematically; no regularity
of the obstacle or of the contact is assumed. Their overlap provides
the first density. Minimality provides the second. Relative isoperimetry
then forces perimeter inside the large constant conformal plateau,
contradicting the competitor bound.}
\label{fig:height:two-densities}
\end{figure}


\subsection{The exterior, its full boundary, and its area}

\begin{proof}[Proof of Theorem~\ref{thm:prepared:deformation}]
\smallskip\noindent\textbf{The connected exterior and its curvature region.}

Choose $F_N$ as in Lemma~\ref{lem:height:BV-existence}.  Its frontier
is compact and, by Lemma~\ref{lem:height:detachment}, has positive
distance from $O_N$.  Every sufficiently small variation of the set
near a frontier point remains admissible.  Thus its boundary is locally
perimeter minimizing in the smooth metric $\widetilde g_N$.

We use the regularity theorem for codimension-one minimizing boundaries:
the support of the perimeter measure is a smooth embedded minimal
hypersurface outside a relatively closed set of Hausdorff dimension at
most $n-8$. A directly applicable smooth-ambient formulation is
\cite[Chapter~7, Section~5, Theorem~5.8]{SimonGMT}; its hypersurface
dimension is $k=n-1$, so the singular bound is $k-7=n-8$.
Multiplicity one for the boundary of a set is recorded in
\cite[Chapter~7, Section~5, Remark~5.2]{SimonGMT}. The compactness and
tangent-cone conclusions used in the minimizing-boundary regularity
argument are Theorems~5.3 and~5.5 of that section.

We also specify the representative of the filled set. Since the
frontier is detached from the obstacle, both insertion and deletion
of a sufficiently small ball are admissible. The coarea and
isoperimetric argument used in the preceding lemma, now applied to
each side, gives positive lower volume densities for both $F_N$ and
its complement at every point in the support of perimeter. On a ball
not meeting that support the indicator has zero distributional
derivative, hence is constant almost everywhere. Let $U_N$ be the
union of those balls on which that constant is zero, and let the
closed filled representative be $\mathcal M_N\setminus U_N$.
The two density bounds identify its topological frontier with the
perimeter support: every neighborhood of a support point contains
both filled and unfilled points. Conversely a point outside the
support has a constant-indicator neighborhood. This representative
agrees with $F_N$ up to a null set; the regularity theorem makes its
frontier $n$-dimensional-volume null. Its perimeter and all the
preceding minimizing comparisons are therefore unchanged.

Put $\Sigma_N=\partial U_N$.
There is only one component of $U_N$.  Indeed there is exactly one
unbounded component, since $F_N$ is bounded and the ambient manifold
has just one end.  Filling all bounded complementary components cannot
increase perimeter.  This statement can be seen directly on regular
boundary charts, where it deletes the separating sheets and adds no
sheet; the singular set has zero $\mathcal H^k$ measure.  Equivalently
it is the component decomposition and perimeter additivity for
finite-perimeter sets \cite{AFP}.  The resulting set is still an
admissible bounded competitor.  If any bounded open component existed,
its positive volume would strictly increase the filled volume; minimality
would either be contradicted by a perimeter decrease or preserved with
a volume increase.  Both alternatives contradict the choice of $F_N$.
Thus $U_N$ is connected.  We take $X_N=\overline{U_N}$, with its
ambient boundary included.

The exterior lies in $\Omega\cap\{|h|<b_3\}$, because the filled
set contains $O_N\supset C_N^h$.  In fact it has a smooth ambient
neighborhood in that region: detachment gives positive separation
between $\Sigma_N$ and $O_N$, and compactness permits a smooth compact
domain containing $O_N$ whose boundary is still strictly inside the
filled set.  Remove that domain.  On the remaining smooth ambient
exterior, Lemma~\ref{lem:height:curvature} gives nonnegative scalar
curvature, also on a neighborhood of $\Sigma_N$.

All local variations of $\Sigma_N$ have zero first variation.  More
precisely, if $Y$ is a compactly supported smooth ambient vector field
near the frontier and $\Phi_s$ is its flow, then $\Phi_s(F_N)$ remains
admissible for both signs of small $s$.  Differentiation of the perimeter
area formula gives
\begin{equation}\label{eq:height:weak-minimality}
 \int_{\partial^*F_N}\divg_{T_x\partial^*F_N}Y\,
 d\mathcal H^k_{\widetilde g_N}=0.
\end{equation}
This is stationarity of the full boundary varifold, including the
assertion that there is no additional first-variation term at the
singular set.  Global outer minimization follows because every bounded
superset of $F_N$ also contains $O_N$.  The maximal choice of $F_N$
also makes it the outermost member among enclosures with this same
minimum perimeter, although outermostness is not required below.

\smallskip\noindent\textbf{Comparison with full original cuts.}\par\nopagebreak[4]

We verify the comparison with original full cuts in some detail.
Detachment and Lemma~\ref{lem:height:ramp} put $\Sigma_N$ in
$\operatorname{int}\Omega$, separated from a neighborhood of the
entire original filling, including every component of $S$.  There are
bounded smooth open sets $F_j$ in the filled manifold such that
\begin{equation}\label{eq:height:strict-approximation}
 \mathbf1_{F_j}\longrightarrow\mathbf1_{F_N}\text{ in }L^1,
 \qquad P_g(F_j)\longrightarrow P_g(F_N),
\end{equation}
and each $F_j$ still contains a fixed neighborhood of the whole filling.
Here is a way to impose the last condition in strict BV approximation.
Choose a smaller compact neighborhood on which $\mathbf1_{F_N}=1$
almost everywhere, separated from $\Sigma_N$, and a compact exterior
region where it is zero.  Smooth the characteristic function in finitely
many coordinate charts only in the remaining region, using a partition
of unity and leaving those two constant regions fixed.  The resulting
smooth functions $u_j\in[0,1]$ converge in $L^1$ with total variation
converging to $P_g(F_N)$.  Choose $\delta_j\downarrow0$ so slowly that
$\|u_j-\mathbf1_{F_N}\|_{L^1}/\delta_j\to0$.  By coarea there is a
regular value $t_j\in(\delta_j,1-\delta_j)$ for which
\[
 P_g(\{u_j>t_j\})
 \le\frac{\int|du_j|_g\,dV_g}{1-2\delta_j}+o(1)
 =P_g(F_N)+o(1).
\]
The $L^1$ difference of this level set from $F_N$ is at most
$\|u_j-\mathbf1_{F_N}\|_{L^1}/\delta_j$, and hence tends to zero.
Lower semicontinuity gives the reverse perimeter liminf.  Thus
$F_j=\{u_j>t_j\}$ satisfies
Equation~\eqref{eq:height:strict-approximation}.  This is the strict
approximation/coarea construction of \cite{AFP}, localized away from
the two constant regions.

For each $j$ retain the component of $\mathcal M_N\setminus\overline{F_j}$
that contains the distant end, and let $D_j$ be its closure.  A compact
smooth boundary has finitely many components; retaining this exterior
component merely selects some of those entire boundary components.
It preserves smoothness and can only decrease the full boundary area.
Because $F_j$ contains a neighborhood of the whole filling,
$D_j\subset\Omega$, its interior lies in $\operatorname{int}\Omega$,
and $D_j$ is closed and connected with a smooth compact intrinsic
boundary $\Gamma_j$.  It is an admissible full cut in
Definition~\ref{def:fullcuts}.  Thus
\[
 a\le\Area_g(\Gamma_j)\le P_g(F_j).
\]
Taking the limit gives $P_g(F_N)\ge a$.  No component of a cut has
been discarded in evaluating its intrinsic boundary; the passage to
$D_j$ is an admissible choice of the exterior domain itself.  There is
no coincidence with $S$ in this approximation because of the fixed
neighborhood, and the original definition still allows coincidence for
its other competitors.

Since $\widetilde g_N\ge e^{-2\epsilon}g$ on $\Omega$, comparison
on every approximate tangent $k$-plane now yields
\begin{equation}\label{eq:height:original-area}
 \mathcal H^k_{\widetilde g_N}(\Sigma_N)
 =P_{\widetilde g_N}(F_N)
 \ge e^{-k\epsilon}P_g(F_N)\ge e^{-k\epsilon}a.
\end{equation}
The first equality uses the regularity theorem and the zero
$\mathcal H^k$ measure of the singular set.

\smallskip\noindent\textbf{Completeness and the end.}

The filled metric is complete, since
$\widetilde g_N\ge e^{-2\epsilon}g$ and the background filled metric
is complete.  Its restriction to $X_N$ is complete with the frontier
included.  Here completeness is understood in the enclosing-set
convention of the Riemannian input: the intrinsic length distance on
$X_N$ may be extended-valued, and its restriction to each finite-distance
equivalence class is complete.  For the intrinsic-length formulation, a Cauchy sequence
has a subsequence whose successive intrinsic distances are less than
$2^{-j}$.  Join successive terms by curves of length less than
$2^{1-j}$ in $X_N$.  Their concatenation has an ambient endpoint,
by ambient completeness, and that endpoint belongs to $X_N$ because
$X_N$ is closed.  Appending this endpoint gives a curve in $X_N$
whose tail length tends to zero, proving intrinsic convergence of the
subsequence and then of the original Cauchy sequence.  The frontier
is included in the closed-enclosing-set sense used in
the Riemannian input in Section~\ref{sec:comparison}.

Finally the end satisfies $\widetilde g_N-I=O_j(r^{2-n})$.
The conformal law applied to Equation~\eqref{eq:height:end-input}
shows
\[
 R_{\widetilde g_N}=O(r^{-n-\beta})+O(r^{-2n+2}).
\]
If needed decrease $\beta$ so that $0<\beta<n-2$; then the displayed
bound is $O(r^{-n-\beta})$.  All graph errors decay exponentially,
and $\psi_N$ is zero on the end.  Scalar curvature is integrable there
and, being smooth and bounded on the compact ambient neighborhood of
the frontier, integrable on the remainder of $X_N$.
The energy estimate is Lemma~\ref{lem:height:curvature}.  This proves
all the stated properties.
\end{proof}

\begin{remark}[Flux and weak boundary conventions]\label{rem:height:flux}
The scalar flux in Proposition~\ref{prop:scalar:flux} is computed on
the complete smooth filled manifold before an enclosure is chosen:
\[
 \mathfrak F_N=\lim_{R\to\infty}\int_{S_R}g(V,\nu_g)\,dA_g
 =\int_{\mathcal M_N}u\Xi\,dV_g,
 \qquad E(\widehat g)-E=-\frac{\mathfrak F_N}{k\omega}.
\]
Thus no elliptic boundary condition or unaccounted singular-boundary
flux is involved.  If one restricts a smooth vector field to $X_N$,
the Gauss--Green theorem for finite-perimeter sets gives its ordinary
reduced-boundary flux.  For instance at a regular truncation radius,
\[
 \int_{U_N\cap B_R}\divg_{\widetilde g_N}Y\,dV_{\widetilde g_N}
 =\int_{S_R}\widetilde g_N(Y,\nu_R)\,dA_{\widetilde g_N}
 -\int_{\partial^*F_N}\widetilde g_N(Y,\nu_{F_N})\,
 dA_{\widetilde g_N},
\]
when $Y$ is supported in the indicated truncation or has its stated
outer trace.  The normal $\nu_{F_N}$ points from the filled set toward
the end.  There is no further measure on the singular set.  This flux
convention and Equation~\eqref{eq:height:weak-minimality} are the weak
identities attached to the constructed boundary.
\end{remark}
